\section{Introduction}
\label{sec:introduction}

Scientific software is often assembled from several repositories, each written to follow the
definitions of a field outside computer science, and the people who rely on it may not accept
a rewrite. The assignment asks for a narrower change. Find one component in which repeated
expensive calculations can be replaced by stored results, show with unit tests that the
results do not change, and only then add parallelism.

The program we were given~\cite{russell2026assignment} works over the finite field $\Fq$,
using the library \code{algebra-finite-fields}~\cite{russell2026fields}. It fills two $N \times N$
matrices $A$ and $B$ with random field elements, computes $C = AB$, computes $\det(A)$,
$\det(B)$ and $\det(C)$ independently, and checks that $\det(C) = \det(A)\det(B)$. The
assignment text reports 6~h~40~min for $N = 4000$ on one core of an Apple M3 Max, and
20~min~16~s for a cached version on the same machine.

This report covers five pieces of work.
\begin{itemize}
  \item We traced the dependency chain and profiled the original program to find where the
        time goes (Sections~\ref{sec:background} and~\ref{sec:profiling}).
  \item We wrote a cached finite field that can be dropped in place of the original
        (Section~\ref{sec:cached}).
  \item We wrote unit tests that compare the cached field with the original on every
        element and every pair of elements (Section~\ref{sec:equivalence}).
  \item We measured run time and peak memory before and after
        (Section~\ref{sec:sequential}).
  \item We wrote an MPI program for the determinant and measured where it starts to beat the
        cached sequential program (Sections~\ref{sec:mpidesign} and~\ref{sec:mpiresults}).
\end{itemize}
